% ============================================================
\newpage
\section{Parte C — Problemas tipo examen}

\begin{solucion}{C1 · Diagnóstico de un servidor}
\textbf{(a)} $P(A)P(U)P(C\mid A,U)P(M\mid C)P(Q\mid C)$: $1+1+4+2+2=\mathbf{10}$ parámetros frente a $2^5-1=31$.

\textbf{(b)}
$A\indep U$: \textbf{sí} ($A\to C\leftarrow U$, colisionador sin evidencia).
$A\indep U\mid M$: \textbf{no} ($M$ es descendiente de $C$: abre el colisionador).
$M\indep Q$: \textbf{no} ($M\leftarrow C\to Q$, causa común no observada).
$M\indep Q\mid C$: \textbf{sí} (causa común observada).
$A\indep Q\mid C$: \textbf{sí} (cadena $A\to C\to Q$ bloqueada).

\textbf{(c)} Factores tras restringir: $P(A)$, $P(U)$, $P(C\mid A,U)$, $P(+m\mid C)=(0.9,\,0.1)$. $Q$ es hoja no observada: se descarta.

\emph{Eliminar $C$}: $f_1(A,U)=\sum_cP(c\mid A,U)P(+m\mid c)$:
\begin{center}\small
\begin{tabular}{cc|l}
$A$ & $U$ & $f_1(A,U)$\\\hline
$+a$ & $+u$ & $0.95(0.9)+0.05(0.1)=0.86$\\
$+a$ & $-u$ & $0.80(0.9)+0.20(0.1)=0.74$\\
$-a$ & $+u$ & $0.60(0.9)+0.40(0.1)=0.58$\\
$-a$ & $-u$ & $0.05(0.9)+0.95(0.1)=0.14$
\end{tabular}
\end{center}
\emph{Eliminar $U$}: $f_2(A)=\sum_uP(u)f_1(A,u)$: $f_2(+a)=0.1(0.86)+0.9(0.74)=0.752$; $f_2(-a)=0.1(0.58)+0.9(0.14)=0.184$.

\emph{Unir con $P(A)$ y normalizar}: $0.05(0.752)=0.0376$; $0.95(0.184)=0.1748$; $P(+m)=0.2124$.
\[P(+a\mid+m)=\frac{0.0376}{0.2124}=\mathbf{0.1770}.\]

\textbf{(d)} Con $+u$ no se suma $U$: pesos $P(A)P(+u)f_1(A,+u)$: $0.05(0.1)(0.86)=0.0043$ y $0.95(0.1)(0.58)=0.0551$.
$P(+a\mid+m,+u)=0.0043/0.0594=\mathbf{0.0724}$. La actualización defectuosa ya explica la CPU alta y la alerta: la creencia de ataque cae de
17.7\,\% a 7.2\,\% (\emph{explaining away}; recuerda que $A\indep U$ sin evidencia, pero no dada $M$).

\textbf{(e)} Ahora cada término incluye $P(+q\mid c)$: $g_1(A,U)=\sum_cP(c\mid A,U)P(+m\mid c)P(+q\mid c)$, con $P(+m\mid c)P(+q\mid c)=0.63$ ($+c$) y $0.005$ ($-c$):
$g_1=0.59875,\ 0.505,\ 0.38,\ 0.03625$ para $(+a,+u),(+a,-u),(-a,+u),(-a,-u)$.
$g_2(+a)=0.1(0.59875)+0.9(0.505)=0.514375$; $g_2(-a)=0.1(0.38)+0.9(0.03625)=0.070625$.
Pesos $0.05(0.514375)=0.02572$ y $0.95(0.070625)=0.06709$: $P(+a\mid+m,+q)=\mathbf{0.2771}$.
$M$ y $Q$ solo informan sobre $C$. Aun sabiendo con certeza $+c$, $P(+a\mid+c)=0.2900$: ese es el techo; la segunda señal acerca $C$ a la certeza
pero no distingue ataque de actualización.
\end{solucion}

\begin{solucion}{C2 · Enumeración frente a eliminación}
\textbf{(a)} Mundos $P(d)P(i)P(c\mid d,i)P(+e\mid i)P(+r\mid c)$:
\begin{center}\small
\begin{tabular}{ccc|l@{\qquad}ccc|l}
$I$&$D$&$C$& & $I$&$D$&$C$& \\\hline
$+i$&$+d$&$+c$& $0.4(0.3)(0.7)(0.8)(0.9)=0.06048$ & $-i$&$+d$&$+c$& $0.4(0.7)(0.1)(0.1)(0.9)=0.00252$\\
$+i$&$+d$&$-c$& $0.4(0.3)(0.3)(0.8)(0.2)=0.00576$ & $-i$&$+d$&$-c$& $0.4(0.7)(0.9)(0.1)(0.2)=0.00504$\\
$+i$&$-d$&$+c$& $0.6(0.3)(0.9)(0.8)(0.9)=0.11664$ & $-i$&$-d$&$+c$& $0.6(0.7)(0.4)(0.1)(0.9)=0.01512$\\
$+i$&$-d$&$-c$& $0.6(0.3)(0.1)(0.8)(0.2)=0.00288$ & $-i$&$-d$&$-c$& $0.6(0.7)(0.6)(0.1)(0.2)=0.00504$
\end{tabular}
\end{center}
Sumas: $0.18576$ y $0.02772$; $P(+r,+e)=0.21348$; $P(+i\mid+r,+e)=\mathbf{0.8702}$.

\textbf{(b)} Los factores $f_1(C,I)$ y $f_2(I)=(0.774,\,0.396)$ de B8 no cambian; solo se añade $P(+e\mid I)=(0.8,\,0.1)$ al final:
$0.3(0.8)(0.774)=0.18576$ y $0.7(0.1)(0.396)=0.02772$ $\Rightarrow$ $\mathbf{0.8702}$.

\textbf{(c)} Enumeración: $2^2=4$ mundos por valor de $I$, 8 productos de 5 factores (32 multiplicaciones). Eliminación: el mayor producto es
$P(D)P(C\mid D,I)$ con alcance $\{C,D,I\}$: 8 entradas. En redes grandes la diferencia es exponencial.

\textbf{(d)} $P(+i\mid+c)=\frac{0.246}{0.442}=\mathbf{0.5566}$; $P(+i\mid+c,+d)=\frac{0.084}{0.112}=\mathbf{0.75}$; $P(+i\mid+c,-d)=\frac{0.162}{0.330}=\mathbf{0.4909}$.
Un curso difícil hace \emph{menos} probable una nota alta; si aun así la nota fue alta, hace falta una explicación más fuerte: inteligencia.
No contradice el \emph{explaining away}: ambos son razonamiento intercausal a través del colisionador $C$ observado. Cuando las dos causas
\emph{favorecen} el efecto, confirmar una baja la otra; cuando una lo \emph{desfavorece} (como $D$), confirmarla sube la otra.
\end{solucion}

\begin{solucion}{C3 · Comparar métodos de muestreo}
\textbf{(a)} Con $+m$: muestras 2 $(-a,+u,+c)$, 4 $(+a,-u,+c)$, 6 $(-a,-u,-c)$, 12 $(-a,+u,+c)$. Se aceptan \textbf{4} y una tiene $+a$: $\hat P=\mathbf{0.25}$
(exacto 0.1770). Esperadas: $12\,P(+m)=12(0.2124)\approx\mathbf{2.5}$; aquí tocaron más.

\textbf{(b)} Peso $=P(+m\mid c)$: $0.9$ si $+c$, $0.1$ si $-c$.
Pesos: $0.1,0.9,\mathbf{0.9},0.1,0.9,0.1,0.1,0.1$ (en negrita la única con $+a$). Suma $3.2$; $\hat P(+a\mid+m)=0.9/3.2=\mathbf{0.2813}$.

\textbf{(c)} $P(A\mid+u,+c)\propto P(A)P(+c\mid A,+u)$: $+a$: $0.05(0.95)=0.0475$; $-a$: $0.95(0.6)=0.57$; $P(+a\mid\cdot)=\mathbf{0.0769}$.
Con $A=-a$: $P(C\mid-a,+u,+m)\propto P(C\mid-a,+u)P(+m\mid C)$: $+c$: $0.6(0.9)=0.54$; $-c$: $0.4(0.1)=0.04$; $P(+c\mid\cdot)=\mathbf{0.9310}$.

\textbf{(d)} Rechazo descartó 8 de 12 muestras (desperdicio $\approx1-P(+m)$, es decir 79\,\%). Verosimilitud con pesos usa todas, pero con pesos desiguales.
Gibbs no desperdicia muestras, pero están correlacionadas. Con evidencia $+m,+q$ ($P=0.0928$, y menor aún con más evidencia) rechazo es el peor;
preferiría verosimilitud con pesos (la evidencia es hoja: pesos no tan extremos) o Gibbs.
\end{solucion}

\begin{solucion}{C4 · Clasificador bayesiano ingenuo}
\textbf{(a)} $J\to\text{Cielo}$, $J\to\text{Viento}$: los atributos son independientes dada la clase.
$\hat j=\argmax_jP(j)P(\text{cielo}\mid j)P(\text{viento}\mid j)$.

\textbf{(b)} $P(J=\text{sí})=6/10$. Clase sí (6): cielo sol 2, nub 3, lluvia 1; viento sí 1, no 5. Clase no (4): sol 2, nub 0, lluvia 2; viento sí 4, no 0.
\begin{center}\small
\begin{tabular}{c|ccc|cc}
 & sol & nub & lluvia & viento sí & viento no\\\hline
$J=$sí & 2/6 & 3/6 & 1/6 & 1/6 & 5/6\\
$J=$no & 2/4 & 0 & 2/4 & 4/4 & 0
\end{tabular}
\end{center}

\textbf{(c)} (lluvia, sí): sí $0.6\cdot\frac16\cdot\frac16=0.0167$; no $0.4\cdot\frac12\cdot1=0.2$ $\Rightarrow P(\text{sí})=0.0769$: \textbf{no}.
(nub, sí): sí $0.6\cdot\frac12\cdot\frac16=0.05$; no $0.4\cdot0\cdot1=0$ $\Rightarrow P(\text{sí})=1$: \textbf{sí}.
(sol, no): sí $0.6\cdot\frac13\cdot\frac56=0.1667$; no $0.4\cdot\frac12\cdot0=0$ $\Rightarrow1$: \textbf{sí}.
Problema: un solo conteo cero anula todo el producto y da certeza absoluta a partir de 4 ejemplos.

\textbf{(d)} Laplace (cielo $k=3$, viento $k=2$):
\begin{center}\small
\begin{tabular}{c|ccc|cc}
 & sol & nub & lluvia & viento sí & viento no\\\hline
$J=$sí & 3/9 & 4/9 & 2/9 & 2/8 & 6/8\\
$J=$no & 3/7 & 1/7 & 3/7 & 5/6 & 1/6
\end{tabular}
\end{center}
(lluvia, sí): $0.6\cdot\frac29\cdot\frac28=0.0333$ vs $0.4\cdot\frac37\cdot\frac56=0.1429$ $\Rightarrow P(\text{sí})=\mathbf{0.1892}$ (no).
(nub, sí): $0.0667$ vs $0.0476$ $\Rightarrow\mathbf{0.5833}$ (sí).
(sol, no): $0.15$ vs $0.0286$ $\Rightarrow\mathbf{0.84}$ (sí).
Las decisiones no cambian; cambia la confianza: ya no hay certezas de 1.
\end{solucion}

\begin{solucion}{C5 · EM con dos monedas}
\textbf{(a)} $Z\in\{A,B\}$ oculta con $P(Z)=\frac12$; $Z\to H$ (caras en 10). $L(\theta)=\prod_j\big[\tfrac12\theta_A^{h_j}(1-\theta_A)^{10-h_j}+\tfrac12\theta_B^{h_j}(1-\theta_B)^{10-h_j}\big]$
(el coeficiente binomial es común y se omite).

\textbf{(b)} Ejemplo $h=5$: $0.6^5\,0.4^5=7.96\times10^{-4}$, $0.5^{10}=9.77\times10^{-4}$, $w=0.4491$.
\begin{center}\small
\begin{tabular}{c|ccccc}
$h_j$ & 5 & 9 & 8 & 4 & 7\\\hline
$w_j=P(A\mid h_j)$ & 0.4491 & 0.8050 & 0.7335 & 0.3522 & 0.6472\\
$1-w_j$ & 0.5509 & 0.1950 & 0.2665 & 0.6478 & 0.3528
\end{tabular}
\end{center}

\textbf{(c)} Caras esperadas de $A$: $\sum w_jh_j=21.30$; cruces: $\sum w_j(10-h_j)=8.57$. De $B$: $33-21.30=11.70$ caras y $17-8.57=8.43$ cruces.
$\theta_A=21.30/29.87=\mathbf{0.7130}$; $\theta_B=11.70/20.13=\mathbf{0.5813}$.
La log-verosimilitud pasa de $-33.09$ a $-31.86$. (Una segunda iteración da $\theta_A=0.7453$, $\theta_B=0.5693$.)

\textbf{(d)} Con la moneda conocida: MLE directo, $\theta_A=$ caras de $A$ / lanzamientos de $A$. Intercambiar las etiquetas da la misma
verosimilitud: la solución es simétrica (no identificable); la inicialización elige cuál es «$A$».
\end{solucion}

\begin{solucion}{C6 · Regresión polinomial y generalización}
\textbf{(a)} $X^\mathsf TX=\begin{pmatrix}5&0\\0&10\end{pmatrix}$, $X^\mathsf T\mathbf y=(10,\,-0.5)$; $\hat\theta=(\mathbf{2},\mathbf{-0.05})$.
Predicciones $2.1,2.05,2,1.95,1.9$; residuos $2.1,-1.25,-1.9,-0.85,1.9$; ECM$_{\text{tr}}=13.915/5=\mathbf{2.783}$.
Prueba: $\hat y(3)=1.85$, $\hat y(-3)=2.15$; ECM$_{\text{te}}=(7.25^2+6.65^2)/2=\mathbf{48.39}$.

\textbf{(b)} Con $\phi=(1,x,x^2)$: $\sum x=\sum x^3=0$, $\sum x^2=10$, $\sum x^4=34$, $\sum x^2y=33.9$:
\[X^\mathsf TX=\begin{pmatrix}5&0&10\\0&10&0\\10&0&34\end{pmatrix},\quad X^\mathsf T\mathbf y=\begin{pmatrix}10\\-0.5\\33.9\end{pmatrix}.\]
La fila central da $\theta_1=-0.05$. Las otras: $5\theta_0+10\theta_2=10$, $10\theta_0+34\theta_2=33.9$; $\det=70$:
$\theta_0=\frac{340-339}{70}=\mathbf{0.0143}$, $\theta_2=\frac{169.5-100}{70}=\mathbf{0.9929}$.
Residuos de entrenamiento $0.114,-0.257,0.086,0.143,-0.086$: ECM$_{\text{tr}}=\mathbf{0.0229}$.
Prueba: $\hat y(3)=8.80$, $\hat y(-3)=9.10$; errores $\pm0.3$: ECM$_{\text{te}}=\mathbf{0.09}$.

\textbf{(c)} En los 5 puntos el polinomio reproduce $y$ exactamente (p.\,ej.\ $x=1$: $0.1+0.2333+0.8083-0.0833+0.0417=1.1$): ECM$_{\text{tr}}=0$.
$\hat y(3)=0.1+0.7+7.275-2.25+3.375=9.2$; $\hat y(-3)=0.1-0.7+7.275+2.25+3.375=12.3$; ECM$_{\text{te}}=(0.01+12.25)/2=\mathbf{6.13}$.

\textbf{(d)} Entrenamiento: grado 4 (0) $<$ grado 2 (0.023) $<$ grado 1 (2.78). Prueba: grado 2 (0.09) $<$ grado 4 (6.13) $<$ grado 1 (48.4).
Grado 1: \textbf{subajuste}; grado 4: \textbf{sobreajuste} (sigue el ruido con 5 parámetros y 5 datos); grado 2 coincide con la estructura real.
\end{solucion}

\begin{solucion}{C7 · $k$-means y escala}
\textbf{(a)} Crudo, $\mu_1=(10,1)$, $\mu_2=(12,9)$. Iteración 1: $d^2(C_3)=(401,\,373)$, $d^2(C_4)=(533,\,401)$ $\Rightarrow$ $\{C_1\}$, $\{C_2,C_3,C_4\}$;
$\mu_2=(24.67,6.33)$, $J=271.3$. Iteración 2: $C_2$: $d^2=(68,\,167.6)$ pasa al grupo 1; $C_3$: $(401,\,47.2)$; $C_4$: $(533,\,56.6)$ $\Rightarrow$
$\{C_1,C_2\}$, $\{C_3,C_4\}$; $\mu_1=(11,5)$, $\mu_2=(31,5)$, $J=17+17+10+10=\mathbf{54}$. Iteración 3 sin cambios. \textbf{Agrupa por ingreso.}

\textbf{(b)} Escala: $C_1(1,1)$, $C_2(1.2,9)$, $C_3(3,2)$, $C_4(3.2,8)$. $C_3$: $d^2=(5,\,52.24)$; $C_4$: $(53.84,\,5)$ $\Rightarrow$ $\{C_1,C_3\}$, $\{C_2,C_4\}$;
$\mu_1=(2,1.5)$, $\mu_2=(2.2,8.5)$, $J=4(1.25)=\mathbf{5}$; estable. \textbf{Agrupa por visitas.} (La estandarización $z$ da la misma partición.)

\textbf{(c)} Crudo: $a(C_1)=d(C_1,C_2)=8.246$, $b=\frac{20.025+23.087}2=21.556$, $s=\mathbf{0.617}$.
Escalado: $a=d(C_1,C_3)=2.236$, $b=\frac{8.002+7.338}2=7.670$, $s=\mathbf{0.708}$.

\textbf{(d)} Ninguno es «el correcto» en abstracto: los dos son óptimos locales de criterios distintos. Lo decide el objetivo (¿segmentar por
poder adquisitivo o por frecuencia?) y, por tanto, la elección de escala; la silueta solo compara dentro de una misma escala.
\end{solucion}

\begin{solucion}{C8 · Robot de reciclaje}
\textbf{(a)}
\begin{center}\small
\begin{tabular}{cc|ll}
\toprule
$s$ & $a$ & $s'$ (prob.) & $R$\\\midrule
$H$ & buscar & $H$ (0.7), $L$ (0.3) & $+2$, $+2$\\
$H$ & esperar & $H$ (1) & $+1$\\
$L$ & buscar & $L$ (0.6), $H$ (0.4) & $+2$, $-3$\\
$L$ & esperar & $L$ (1) & $+1$\\
$L$ & recargar & $H$ (1) & $0$\\\bottomrule
\end{tabular}
\end{center}

\textbf{(b)} \emph{Iteración de políticas.}
\begin{itemize}
\item $\pi_0=(\text{esp},\text{esp})$: $V=1+0.9V\Rightarrow V(H)=V(L)=10$. Mejora: $H$: buscar $0.7(2+9)+0.3(2+9)=11>10$; $L$: buscar $0.6(11)+0.4(-3+9)=9$,
esperar 10, recargar 9 $\Rightarrow$ $\pi_1=(\text{buscar},\text{esp})$.
\item $\pi_1$: $V(L)=10$; $V(H)=2+0.63V(H)+0.27(10)\Rightarrow V(H)=4.7/0.37=12.7027$. Mejora: $H$: buscar 12.7027 vs esperar 12.4324;
$L$: buscar $6.6+0.4(-3+11.4324)=9.9730$, esperar 10, recargar $0.9(12.7027)=11.4324$ $\Rightarrow$ $\pi_2=(\text{buscar},\text{recargar})$.
\item $\pi_2$: $V(L)=0.9V(H)$, $V(H)=2+0.63V(H)+0.27(0.9V(H))\Rightarrow0.127V(H)=2$: $V(H)=\mathbf{15.748}$, $V(L)=\mathbf{14.173}$. Mejora: $H$: buscar 15.748
vs esperar 15.173; $L$: buscar 13.323, esperar 13.756, recargar 14.173. Sin cambios: \textbf{$\pi^*=(\text{buscar},\text{recargar})$}.
\end{itemize}

\textbf{(c)} \emph{Iteración de valores.}
\begin{center}\small
\begin{tabular}{c|cc|ccc|c}
\toprule
$k$ & $H$:buscar & $H$:esperar & $L$:buscar & $L$:esperar & $L$:recargar & $(V_k(H),V_k(L))$\\\midrule
1 & 2 & 1 & 0 & 1 & 0 & $(2,1)$\\
2 & 3.53 & 2.8 & 1.26 & 1.9 & 1.8 & $(3.53,1.9)$\\
3 & 4.7369 & 4.177 & 2.2968 & 2.71 & 3.177 & $(4.7369,3.177)$\\
4 & 5.842 & 5.2632 & 3.4209 & 3.8593 & 4.2632 & $(5.842,4.2632)$\\\bottomrule
\end{tabular}
\end{center}
Ej.\ $k=2$, $H$:buscar $=0.7(2+0.9\cdot2)+0.3(2+0.9\cdot1)=3.53$. Con $V_4$ la voraz es (buscar, recargar), la óptima; con $V_1,V_2$ en $L$ ganaba esperar
(horizonte corto). Los valores aún están lejos de $(15.75,14.17)$: convergen geométricamente con razón $0.9$.

\textbf{(d)} Solo cambia $Q(L,\text{buscar})$, que empeora; la política óptima no usa esa acción, así que $\pi^*$ y $V^*$ \textbf{no cambian}.
\end{solucion}

\begin{solucion}{C9 · Q-learning en el corredor}
\textbf{(a)} $Q\leftarrow Q+0.5[y-Q]$; terminales $L,R$ con $y=r$.
\begin{center}\small
\begin{tabular}{c|c|c|c}
Ep. & Transición & $y$ & $Q$ anterior $\to$ nuevo\\\hline
1 & $(2,\text{der},-1,3)$ & $-1+0.9(0)=-1$ & $Q(2,\text{der})$: $0\to-0.5$\\
1 & $(3,\text{der},+10,R)$ & $10$ & $Q(3,\text{der})$: $0\to5$\\
2 & $(1,\text{der},-1,2)$ & $-1+0.9\max(0,-0.5)=-1$ & $Q(1,\text{der})$: $0\to-0.5$\\
2 & $(2,\text{der},-1,3)$ & $-1+0.9(5)=3.5$ & $Q(2,\text{der})$: $-0.5\to1.5$\\
2 & $(3,\text{izq},-1,2)$ & $-1+0.9(1.5)=0.35$ & $Q(3,\text{izq})$: $0\to0.175$\\
2 & $(2,\text{der},-1,3)$ & $-1+0.9(5)=3.5$ & $Q(2,\text{der})$: $1.5\to2.5$\\
2 & $(3,\text{der},+10,R)$ & $10$ & $Q(3,\text{der})$: $5\to7.5$\\
3 & $(2,\text{izq},-1,1)$ & $-1+0.9\max(0,-0.5)=-1$ & $Q(2,\text{izq})$: $0\to-0.5$\\
3 & $(1,\text{izq},+2,L)$ & $2$ & $Q(1,\text{izq})$: $0\to1$\\
4 & $(1,\text{der},-1,2)$ & $-1+0.9(2.5)=1.25$ & $Q(1,\text{der})$: $-0.5\to0.375$\\
4 & $(2,\text{der},-1,3)$ & $-1+0.9(7.5)=5.75$ & $Q(2,\text{der})$: $2.5\to4.125$\\
4 & $(3,\text{der},+10,R)$ & $10$ & $Q(3,\text{der})$: $7.5\to8.75$
\end{tabular}
\end{center}
Tabla final: $Q(1,\cdot)=(1,\ 0.375)$, $Q(2,\cdot)=(-0.5,\ 4.125)$, $Q(3,\cdot)=(0.175,\ 8.75)$ (izq, der).

\textbf{(b)} Voraz: 1 $\to$ \textbf{izq}, 2 $\to$ der, 3 $\to$ der.

\textbf{(c)} $Q^*$ (de B24): $(2.504,\,4.016)$, $(3.495,\,6.134)$, $(5.616,\,8.904)$. El estado 1 está equivocado: $Q(1,\text{der})$ se actualizó cuando
$Q(2,\cdot)$ aún era bajo, y $Q(1,\text{izq})$ recibió de golpe medio $+2$. Hace falta más exploración (y más episodios) para que el valor de la
derecha se propague hasta el estado 1.

\textbf{(d)} El objetivo de Q-learning usa $\max_{a'}Q(s',a')$, no la acción que el agente tomará después: la exploración solo decide qué
entrada se actualiza, sin contaminar los objetivos de las demás. Con SARSA, la transición $(2,\text{der},-1,3)$ seguida de «izq» habría usado
$Q(3,\text{izq})=0$: $y=-1$ y $Q(2,\text{der})=-0.75$ en vez de $1.5$. Si hubiera pedido «der» y resbalado a 2, la tupla sería
$(3,\text{der},-1,2)$ y se actualizaría $Q(3,\text{der})$: el ruido del entorno queda promediado dentro de ese valor.
\end{solucion}

\begin{solucion}{C10 · Monte Carlo frente a TD}
\textbf{(a)} $V(A)$: un solo retorno, $0+0=0$ $\Rightarrow$ $\mathbf0$. $V(B)$: retornos $0,1,1,1,1,1,1,0$ $\Rightarrow$ $6/8=\mathbf{0.75}$.
\textbf{(b)} En el punto fijo batch, cada valor es el promedio de sus objetivos. $B$ siempre termina: $V(B)=0.75$. La única transición de $A$
va a $B$ con 0: $V(A)=0+V(B)=\mathbf{0.75}$.
\textbf{(c)} TD es la respuesta del \emph{modelo de máxima verosimilitud} (A siempre pasa a B con 0, y B paga 1 con prob.\ 3/4) suponiendo Markov;
MC minimiza el error cuadrático sobre los retornos observados (y da 0 por un único episodio). Si el proceso es Markov, TD predice mejor el futuro; MC
no supone Markov.
\textbf{(d)} Ep.\,1: $V(A)\leftarrow0+0.5(0+0-0)=0$; $V(B)\leftarrow0$. Luego seis veces $B\to1$: $0.5,\,0.75,\,0.875,\,0.9375,\,0.96875,\,0.984375$; el último
$B\to0$: $0.4922$. Resultado $V(A)=\mathbf0$, $V(B)=\mathbf{0.4922}$: en línea con $\alpha$ fija, pesan más las experiencias recientes.
\end{solucion}

\begin{solucion}{C11 · Retropropagación con softmax}
\textbf{(a)} $\mathbf z_1=W_1\mathbf x+\mathbf b_1=(2-1+0,\ 1+1-3)=(1,-1)$; $\mathbf h=(1,0)$; $\mathbf z_2=W_2\mathbf h=(1,\,0.5)$;
$\mathbf p=\mathrm{softmax}(1,0.5)=(0.6225,\,0.3775)$; $\ell=-\ln0.3775=\mathbf{0.9741}$.

\textbf{(b)} $\boldsymbol\delta_2=\mathbf p-\mathbf t=(0.6225,-0.6225)$;
$\frac{\partial\ell}{\partial W_2}=\boldsymbol\delta_2\mathbf h^\mathsf T=\begin{pmatrix}0.6225&0\\-0.6225&0\end{pmatrix}$, $\frac{\partial\ell}{\partial\mathbf b_2}=\boldsymbol\delta_2$.
$W_2^\mathsf T\boldsymbol\delta_2=(1(0.6225)+0.5(-0.6225),\ -1(0.6225)+1(-0.6225))=(0.3113,\,-1.2449)$.
$\boldsymbol\delta_1=(0.3113,\,-1.2449)\odot(1,0)=(0.3113,\,0)$;
$\frac{\partial\ell}{\partial W_1}=\boldsymbol\delta_1\mathbf x^\mathsf T=\begin{pmatrix}0.3113&-0.3113\\0&0\end{pmatrix}$, $\frac{\partial\ell}{\partial\mathbf b_1}=(0.3113,0)$.

\textbf{(c)} $\eta=1$: $W_1=\begin{pmatrix}1.6888&1.3112\\1&-1\end{pmatrix}$, $\mathbf b_1=(-0.3113,-3)$, $W_2=\begin{pmatrix}0.3775&-1\\1.1225&1\end{pmatrix}$, $\mathbf b_2=(-0.6225,0.6225)$.
Nuevo forward: $\mathbf z_1=(0.0663,-1)$, $\mathbf h=(0.0663,0)$, $\mathbf z_2=(-0.5974,\,0.6969)$, $\mathbf p=(0.2151,\,0.7849)$, $\ell=\mathbf{0.2422}$.

\textbf{(d)} La segunda unidad oculta está apagada ($z=-1$): no cambian la segunda fila de $W_1$, $b_{1,2}$, ni la segunda columna de $W_2$
(multiplica $h_2=0$). Neurona muerta: si está apagada para \emph{todas} las entradas, nunca recibe gradiente.
\end{solucion}

\begin{solucion}{C12 · Un transformer diminuto}
\textbf{(a)}
\begin{center}\small
\begin{tabular}{l|l|c|c}
Prefijo & Cálculo & $p(\text{correcto})$ & $-\ln p$\\\hline
$\langle s\rangle$ & $e^2/(e^2+3+e^{-1})=7.389/10.757$ & 0.6869 & 0.3756\\
… el & $e^2/(e^2+e+3)=7.389/13.107$ & 0.5637 & 0.5732\\
… gato & $e^2/(e+1+e^2+1+1)$ & 0.5637 & 0.5732\\
… come & $e^3/(1+1+e+1+e^3)=20.09/25.80$ & 0.7784 & 0.2505
\end{tabular}
\end{center}
Pérdida media $=1.7724/4=\mathbf{0.4431}$; perplejidad $e^{0.4431}=\mathbf{1.558}$ (como elegir entre $\approx1.56$ opciones por paso).

\textbf{(b)} $\tau=0.5$ (logits $\times2$): $(0.1137,0.0154,\mathbf{0.8401},0.0154,0.0154)$. $\tau=2$ (logits $/2$): $(0.2238,0.1357,\mathbf{0.3690},0.1357,0.1357)$.

\textbf{(c)} $S=QK^\mathsf T/\sqrt2=\begin{pmatrix}0.7071&0.7071&1.4142\\0&1.4142&1.4142\\1.4142&0&1.4142\end{pmatrix}$. Con máscara:
$A=\begin{pmatrix}1&0&0\\0.1956&0.8044&0\\0.4458&0.1084&0.4458\end{pmatrix}$,
$O=AV=\begin{pmatrix}2&0\\0.3911&1.6089\\1.3374&0.6626\end{pmatrix}$.
(Fila 3: $0.4458(2,0)+0.1084(0,2)+0.4458(1,1)$.)

\textbf{(d)} Embeddings: $50\,000\times512=\mathbf{25\,600\,000}$. $W_Q,W_K,W_V$: $3\times512\times64=98\,304$; $W_O$: $64\times512=32\,768$; atención: $\mathbf{131\,072}$.
El vocabulario domina en modelos pequeños.
\end{solucion}

\begin{solucion}{C13 · Auditoría de un sistema de crédito}
\textbf{(a)} Exactitud global $(48+32+18+68)/200=\mathbf{0.83}$.
\begin{center}\small
\begin{tabular}{c|cccccc}
Grupo & Tasa base & Aprobación & TPR & FPR & PPV & Exactitud\\\hline
A & 0.60 & 0.56 & 0.80 & 0.20 & 0.857 & 0.80\\
B & 0.30 & 0.20 & 0.60 & 0.029 & 0.90 & 0.86
\end{tabular}
\end{center}
\textbf{(b)} Se violan paridad ($0.56$ vs $0.20$), igualdad de oportunidades ($0.80$ vs $0.60$) y \emph{equalized odds}. En B, el 40\,\% de quienes pagarían
es rechazado (costo para esas personas); en A, el 20\,\% de quienes no pagarían es aprobado (costo para el prestamista y riesgo de
sobreendeudamiento). La mejor exactitud de B no compensa esa asimetría.
\textbf{(c)} Las tasas base difieren (0.6 vs 0.3): con PPV (calibración) igual y tasas base distintas, no pueden igualarse también FPR y FNR salvo con un
clasificador perfecto. Etiquetas selectivas: solo se observa si paga quien recibió crédito; los rechazados nunca generan la etiqueta, y el
sesgo se perpetúa en el siguiente entrenamiento.
\textbf{(d)} Alto riesgo (evaluación de solvencia de personas físicas, anexo III): gestión de riesgos, gobernanza de datos, documentación técnica,
registros, transparencia, supervisión humana, exactitud y robustez, y evaluación de conformidad antes de ponerlo en el mercado.
\textbf{(e)} $40(0.3)=12$ Wh por solicitud; por solicitud bien resuelta $12/0.8=\mathbf{15\ Wh}$; $10\,000$ solicitudes diarias: $\mathbf{120\ kWh/día}$.
\end{solucion}
